\slajdid{04-s035}
\begin{frame}[label=04-s035]{Ponovljeno komplementiranje}
% element: 04-s035:e01 — Dva prikaza nule kod komplementa maksimuma
% element: 04-s035:e02 — Asimetrija drugog komplementa i veća apsolutna negativna vrednost
% element: 04-s035:e03 — Isto komplementiranje za apsolutnu vrednost; potrebno poznavanje koda
% element: 04-s035:e04 — Obe izvorne formule ponovljenog komplementiranja
\begin{itemize}
\item Komplement maksimalne vrednosti: \alert{dva prikaza nule}.
\item Komplement osnove: izuzimajući nulu, nejednak broj pozitivnih i negativnih vrednosti. Najveća apsolutna negativna vrednost veća je od najvećeg pozitivnog broja.
\end{itemize}
Za negativan broj apsolutnu vrednost dobijamo istim komplementiranjem. Moramo znati način prikazivanja.
{\fontsize{10}{12}\selectfont
\[111\ldots11-\overline{b_{n-1}}\,\overline{b_{n-2}}\,\overline{b_{n-3}}\ldots\overline{b_1}\,\overline{b_0}=b_{n-1}b_{n-2}b_{n-3}\ldots b_1b_0\]
\[111\ldots11-\left(\overline{b_{n-1}}\,\overline{b_{n-2}}\,\overline{b_{n-3}}\ldots\overline{b_1}\,\overline{b_0}+1\right)+1=b_{n-1}b_{n-2}b_{n-3}\ldots b_1b_0\]}

\end{frame}
